% ============================================================
%  Pre-registro del experimento 2 - TFG AutoPipe
%  Autor: Ricardo Antonio Borbon Mena
%
%  POR QUE ES UN DOCUMENTO SEPARADO
%  --------------------------------
%  El pre-registro v1.0 declara en su primera pagina que no
%  modifica ningun criterio de exito del anteproyecto. El
%  experimento 2 si toca parametros que la seccion 4.3 fija
%  (objetivo, presupuesto, corpus, variable primaria), de modo
%  que no puede ser una version de aquel documento sin
%  contradecirlo. Aquel gobierna el experimento 1 y queda
%  intacto, con su etiqueta prereg-v1.0.
%
%  REGLA DE ORO
%  ------------
%  Ninguna version se escribe DESPUES de la ejecucion que
%  gobierna. Este documento se cierra con un commit propio y
%  una etiqueta anotada (prereg-v2.0) creada el MISMO dia,
%  antes de ejecutar el experimento 2.
%
%  Compilar con:  pdflatex Preregistro-E2.tex
% ============================================================

\documentclass[11pt,letterpaper]{article}

\usepackage[utf8]{inputenc}
\usepackage[T1]{fontenc}
\usepackage[spanish,es-tabla]{babel}
\usepackage{lmodern}
\usepackage[letterpaper,top=2.4cm,bottom=2.4cm,left=2.8cm,right=2.5cm]{geometry}
\usepackage{xcolor}
\usepackage{colortbl}
\usepackage{booktabs}
\usepackage{array}
\usepackage{tabularx}
\usepackage{enumitem}
\usepackage{fancyhdr}
\usepackage{titlesec}
\usepackage{amsmath}

\definecolor{tecnavy}{HTML}{003366}
\definecolor{gris}{HTML}{666666}
\definecolor{tablehead}{HTML}{E6EEF5}
\definecolor{avisofondo}{HTML}{FDF3E7}
\definecolor{avisoborde}{HTML}{C87A28}

\newcolumntype{Y}{>{\raggedright\arraybackslash}X}
\newcolumntype{L}[1]{>{\raggedright\arraybackslash}p{#1}}

\setlength{\parindent}{0pt}
\setlength{\parskip}{5pt}
\renewcommand{\arraystretch}{1.2}
\setlist[itemize]{leftmargin=1.4em,topsep=2pt,itemsep=3pt}
\setlist[enumerate]{leftmargin=1.6em,topsep=2pt,itemsep=3pt}

\titleformat{\section}{\large\bfseries\sffamily\color{tecnavy}}{\thesection}{0.6em}{}
\titleformat{\subsection}{\normalsize\bfseries\sffamily}{\thesubsection}{0.5em}{}

\pagestyle{fancy}
\fancyhf{}
\renewcommand{\headrulewidth}{0.4pt}
\fancyhead[L]{\small\textit{AutoPipe}}
\fancyhead[R]{\small Pre-registro del experimento 2 --- \preregversion}
\fancyfoot[C]{\thepage}

% --- Identidad de la version --------------------------------
\newcommand{\preregversion}{v2.0}
\newcommand{\preregtag}{prereg-v2.0}
\newcommand{\preregfecha}{19 de septiembre de 2026}

% Caja de aviso, para lo que no se puede pasar por alto.
% Se define como comando y no como entorno a propósito: tabularx lee su cuerpo
% entero de una vez y no admite abrirse en la parte inicial de un \newenvironment.
\setlength{\fboxsep}{9pt}
\newcommand{\aviso}[1]{%
  \par\vspace{7pt}\noindent
  \fcolorbox{avisoborde}{avisofondo}{%
    \parbox{\dimexpr\textwidth-2\fboxsep-2\fboxrule\relax}{#1}%
  }\par\vspace{7pt}}

\begin{document}

% =========== PORTADA ========================================
\begin{center}
{\LARGE\bfseries\color{tecnavy} Pre-registro del experimento 2}\\[6pt]
{\large AutoPipe --- Trabajo Final de Graduación}\\[4pt]
{\small Ricardo Antonio Borbón Mena --- Instituto Tecnológico de Costa Rica}\\[14pt]

\begin{tabularx}{0.92\textwidth}{L{3.6cm} Y}
\toprule
\rowcolor{tablehead} \multicolumn{2}{l}{\textbf{Identificación de esta versión}} \\
\midrule
Versión & \preregversion \\
Fecha de cierre & \preregfecha \\
Estado & \textbf{Versión congelada previa a la ejecución del experimento~2.} Se
cierra después de la aprobación del profesor asesor registrada en el ADR-004
(2026-09-19) y \textbf{antes} de ejecutar un solo ensayo: ningún resultado del
experimento~2 existía cuando se escribió, y el arnés que lo ejecuta se escribe
después de este documento. \\
Etiqueta de git & \texttt{\preregtag}, etiqueta anotada que apunta al \emph{commit}
que cierra esta versión y creada el mismo día que él. El fechado verificable es la
fecha de \textbf{ese commit}, y se obtiene con \texttt{git log -1 \preregtag}. \\
Forma de cita en el informe & «Pre-registro del experimento 2, \preregversion,
etiqueta \texttt{\preregtag}, \preregfecha» \\
Gobierna & El \textbf{experimento 2} únicamente. El experimento 1 lo gobierna el
pre-registro v1.0, etiqueta \texttt{prereg-v1.0}, que \textbf{no se modifica}. \\
\bottomrule
\end{tabularx}
\end{center}

\vspace{10pt}

% =========== 1. PROPOSITO ===================================
\section{Propósito y relación con lo ya comprometido}
\label{sec:proposito}

El experimento 1 se ejecutó el 2026-09-13 bajo el pre-registro v1.0 y la hipótesis
no se confirmó. Su análisis dejó al descubierto que la medición tenía poca
resolución, y el diagnóstico registrado en el ADR-004 lo cuantificó. Este documento
fija, antes de ejecutarlo, el protocolo de un \textbf{segundo experimento} que mide
la misma hipótesis con un instrumento capaz de detectarla.

Tres propiedades hacen que esto no sea una modificación retroactiva del primero, y
conviene enunciarlas antes que nada:

\begin{enumerate}
  \item \textbf{El experimento 1 queda intacto.} Sus números, su criterio y su
        veredicto se reportan tal como salieron, junto a los de este. El
        pre-registro v1.0 no se edita.
  \item \textbf{La Sección~4.3 del anteproyecto ya está cumplida.} El experimento
        comprometido ante el tribunal se ejecutó completo. Lo que este documento
        establece es \textbf{adición}, nunca sustitución. Ningún entregable ni
        criterio se retira.
  \item \textbf{Este documento se cierra y se etiqueta antes de ejecutar}, como los
        tres anteriores.
\end{enumerate}

A diferencia del pre-registro v1.0 ---que solo cerraba silencios de §4.3---, este
documento \textbf{sí modifica parámetros que §4.3 fija}: el objetivo, el
presupuesto, el tamaño del corpus y la variable primaria. Por eso es un documento
separado y por eso el ADR-004 exige la aprobación del profesor asesor.

% =========== 2. DECLARACIONES ===============================
\section{Declaraciones obligatorias}
\label{sec:declaraciones}

Dos hechos podrían leerse como grados de libertad ocultos si no se dijeran. Se
dicen aquí, en la primera sección de contenido, y no en una nota al pie.

\aviso{%
\textbf{Primera. El umbral del 99\,\% se eligió midiendo sobre datos ya ejecutados.}

El diagnóstico del ADR-004 recalculó, sobre las curvas del experimento 1, cuántos
ensayos necesita el arranque en frío bajo objetivos del 95, 97, 99, 99,5 y 100\,\%.
La mediana pasa de 3 a 10 ensayos entre el 95 y el 99\,\%, y a partir de ahí no
mejora: por eso se elige el 99\,\%.

Esa medición \textbf{leyó únicamente la condición de arranque en frío}. Ninguna
curva de ninguna condición caliente se consultó, y el guion que la produjo lo impide
por construcción. La distinción es la que separa calibrar el instrumento de ajustar
el resultado: el umbral se fijó por dónde satura el frío, que es una propiedad del
espacio y del corpus, y no por dónde aparece una diferencia entre condiciones.
}

\aviso{%
\textbf{Segunda. Diecisiete de los cuarenta y dos conjuntos se usaron antes.}

El corpus del experimento 2 son 42 conjuntos (Sección~\ref{sec:corpus}), y 17 de
ellos ---los que se incorporan--- se usaron en la semana 8 para la evaluación
exploratoria de las cuatro estrategias del recomendador, junto con los diez que
siguen excluidos por la regla de familia.

Aquella evaluación \textbf{no seleccionó nada}: las cuatro estrategias empataron y
ninguna decisión de diseño se tomó a partir de ella. La estrategia
\texttt{frequent} se añadió por el argumento conceptual de aislar la contribución
del meta-aprendizaje, no por haber ganado. Aun así, esos 17 conjuntos ya fueron
observados, y quien lea los resultados tiene derecho a saberlo sin deducirlo.
}

% =========== 3. LO QUE SE HEREDA ============================
\section{Lo que se hereda del pre-registro v1.0 sin cambios}
\label{sec:herencia}

Todo lo que este documento no menciona se hereda literalmente del pre-registro v1.0
y no se repite aquí. En particular, y por ser lo que más importa:

\begin{itemize}
  \item \textbf{El modelo base congelado}: \texttt{KNeighbors\allowbreak Classifier}
        con $k=5$, pesos uniformes, distancia de Minkowski con $p=2$ y algoritmo de
        fuerza bruta.
  \item \textbf{El espacio de búsqueda \texttt{tabular-v1}}, sin tocar: los mismos
        cuatro pasos en el mismo orden, las mismas opciones, las mismas 144
        configuraciones con la omisión de pasos. No se crea ningún
        \texttt{tabular-v2}.
  \item \textbf{La extracción de meta-características}: biblioteca, grupos,
        submuestreo del grupo de complejidad, imputación previa, determinismo y
        regla de descarte de no finitas.
  \item \textbf{Las cuatro garantías del protocolo de evaluación} y la ausencia de
        fuga en el nivel meta: el recomendador de cada pliegue se ajusta solo con
        los conjuntos restantes.
  \item \textbf{Los ensayos inyectados cuentan}, y los ensayos inevaluables consumen
        su turno.
  \item \textbf{El desempate canónico} entre configuraciones ganadoras empatadas.
\end{itemize}

Que el espacio no cambie no es una omisión: es lo que permite calcular el óptimo
exacto de la Sección~\ref{sec:objetivo}, y lo que mantiene interpretable la base de
meta-conocimiento ya construida.

% =========== 4. EL OBJETIVO =================================
\section{El objetivo: una vara externa e igual para ambas condiciones}
\label{sec:objetivo}

\subsection{Definición}

Para cada pareja (conjunto de datos $d$, semilla $s$) se define el \textbf{óptimo de
referencia} $M(d,s)$ como el mejor valor de macro-F1 que alcanza \emph{alguna} de las
144 configuraciones de \texttt{tabular-v1}, evaluada sobre \textbf{la misma
partición} que usarán las búsquedas de esa pareja: partición estratificada con el
20\,\% de validación y semilla $s$.

El objetivo de esa pareja es
\[
  \mathrm{objetivo}(d,s) \;=\; 0{,}99 \times M(d,s).
\]

$M(d,s)$ se obtiene por \textbf{rejilla exhaustiva}, no por estimación: las 144
configuraciones se evalúan todas. Es factible precisamente porque el espacio es
pequeño, y no lo sería en un espacio con hiperparámetros continuos.

\subsection{Tres consecuencias que hay que enunciar}

\textbf{La vara deja de ser autorreferencial.} En el experimento 1 el objetivo era el
95\,\% del mejor valor que \emph{el propio arranque en frío} había logrado con su
presupuesto. Eso hacía que la condición de control definiera el listón que la
condición experimental debía superar. Ahora ambas se miden contra el mismo número,
que no depende de ninguna de las dos.

\textbf{El arranque en frío ya puede censurarse.} La v1.0 demostraba que el frío
nunca podía quedar censurado, porque el objetivo era una fracción de su propio
máximo. Esa demostración \textbf{ya no vale}, y su desaparición es intencionada: una
condición que no puede fallar tampoco puede informar. En consecuencia, la
comprobación activa que el arnés del experimento 1 hace ---abortar si el frío queda
censurado--- \textbf{debe retirarse} para el experimento 2, y no puede retirarse por
descuido: se retira aquí, por escrito.

\textbf{El regret nunca es negativo.} Como $M(d,s)$ es el máximo sobre el mismo
espacio y la misma partición, ningún ensayo puede superarlo. Es un invariante
comprobable, y se comprobará con una prueba: si alguna curva excediera su referencia,
indicaría que la partición de la rejilla y la de la búsqueda no coinciden.

\subsection{Casos límite}

Si para alguna pareja resultara $M(d,s) \leq 0$, el regret normalizado no está
definido: esa pareja se excluye y se reporta cuántas fueron. No se espera ninguna
---macro-F1 es positivo en todos los conjuntos del corpus--- pero la regla se fija
ahora y no cuando aparezca el caso.

% =========== 5. PRESUPUESTO Y CENSURA =======================
\section{Presupuesto y censura}
\label{sec:presupuesto}

\begin{center}
\begin{tabularx}{\textwidth}{L{5.2cm} L{2.6cm} Y}
\toprule
\rowcolor{tablehead} \textbf{Parámetro} & \textbf{Valor} & \textbf{Por qué} \\
\midrule
Ensayos por condición & \textbf{60} &
Con el objetivo al 99\,\%, la peor corrida del experimento 1 necesitó 30 ensayos
exactos, justo el presupuesto. Mantenerlo en 30 haría que la censura dominara el
análisis. \\
\midrule
Valor de censura & \textbf{61} &
Presupuesto más uno, la misma regla de la v1.0. Aplica solo a la variable
secundaria: el \emph{regret} no se censura. \\
\midrule
Censuras de robustez & \textbf{120 y 600} &
Castigos más severos, en la misma proporción que los 60 y 300 de la v1.0 respecto de
su presupuesto. \\
\bottomrule
\end{tabularx}
\end{center}

% =========== 6. LA VARIABLE PRIMARIA ========================
\section{La variable primaria: \textit{regret} normalizado}
\label{sec:variable}

\subsection{Definición}

Sea una corrida con valores $v_1,\dots,v_T$ ($T=60$), donde los ensayos inevaluables
valen \texttt{NaN}. Sea $b_t$ el mejor valor finito entre los primeros $t$ ensayos,
indefinido si ninguno lo es. Se define el regret instantáneo
\[
  r_t \;=\;
  \begin{cases}
    1 & \text{si } b_t \text{ no está definido},\\[2pt]
    \dfrac{M(d,s) - b_t}{M(d,s)} & \text{en otro caso},
  \end{cases}
\]
y la variable primaria como su promedio sobre el presupuesto:
\[
  R \;=\; \frac{1}{T}\sum_{t=1}^{T} r_t \;\in\; [0,1].
\]

\textbf{Menor es mejor.} Es el área normalizada entre la curva de mejor-hasta-ahora y
el óptimo de referencia.

\subsection{Por qué esta variable y no la del experimento 1}

Los ensayos hasta el objetivo son un entero pequeño. Sobre los 25 conjuntos del
experimento 1 tomó \textbf{7 valores distintos} y produjo \textbf{18 colisiones}. La
prueba de Wilcoxon descarta los pares empatados, de modo que el contraste primario
operó sobre una fracción de la muestra. El mismo cálculo con el regret da
\textbf{25 valores distintos y ninguna colisión}.

El regret además \textbf{usa los 60 ensayos}, no solo el primero que cruza un umbral:
distingue una búsqueda que llega tarde al objetivo pero mejora sostenidamente de otra
que lo cruza pronto y se estanca. Y trata los ensayos inevaluables de forma coherente
con la v1.0: consumen su turno y no reducen el regret.

\subsection{Los ensayos hasta el objetivo se conservan}

La variable del experimento 1 se mide igual y se reporta como \textbf{análisis
secundario declarado}, con la definición del anteproyecto y su criterio del 30\,\%.
Es lo que permite comparar los dos experimentos en los mismos términos y decir con
precisión qué cambió al cambiar el instrumento.

% =========== 7. CORPUS Y SEMILLAS ===========================
\section{Corpus, semillas y condiciones}
\label{sec:corpus}

\subsection{Corpus: 42 conjuntos, no 52}

El cribado de la semana 5 descargó y verificó \textbf{52} candidatos, de los que
seleccionó 25. El experimento 2 amplía el corpus a \textbf{42}, y el número no es 52
por una razón que conviene explicar porque es la que más podría sesgar este
experimento.

De los 27 candidatos que quedaron fuera del corpus del experimento 1, \textbf{diez lo
hicieron por la regla de familia} ---seis de la familia «dígitos manuscritos» y
cuatro de la familia «NASA MDP»--- y diecisiete simplemente quedaron por debajo del
corte en el orden por sensibilidad. \textbf{Los diez excluidos por familia siguen
excluidos.} Entran los diecisiete restantes:
\[
  25 + 17 = 42 .
\]

\textbf{Por qué la regla de familia es más importante aquí que en el experimento 1.}
En una validación por exclusión de conjunto, admitir dos conjuntos de la misma
familia significa que, en el pliegue que excluye a uno, su casi gemelo permanece en
los datos con los que se ajusta el recomendador. El vecino más cercano de la consulta
sería entonces un hermano suyo, y el arranque en caliente parecería acertar por una
similitud que no es la que el método pretende explotar. El sesgo iría
\textbf{a favor} de la hipótesis, que es la dirección en la que ningún sesgo puede
admitirse.

Los criterios de inclusión del anteproyecto ---licencia abierta, 500 a 50\,000
instancias, hasta 100 características, hasta 10 clases, sin imagen ni texto
embebidos--- se mantienen sin alteración: los 42 ya los cumplen. Lo único que se
retira es el \emph{corte} por sensibilidad en el orden por IQR, que existía para
quedarse con 25 y que el anteproyecto imponía por el límite de 25 conjuntos.

Véase la Sección~\ref{sec:declaraciones}, segunda declaración.

\subsection{Semillas: diez}

$\{101,\,202,\,303,\,404,\,505,\,606,\,707,\,808,\,909,\,1010\}$, \textbf{idénticas en
todas las condiciones}. Las cinco primeras son las de la v1.0. Se pasa de cinco a
diez porque el cómputo lo permite y porque la observación que entra en la prueba es
la mediana por conjunto: diez semillas la estabilizan.

\subsection{Las cinco condiciones}

\begin{center}
\begin{tabularx}{\textwidth}{L{2.3cm} L{2.7cm} Y}
\toprule
\rowcolor{tablehead} \textbf{Condición} & \textbf{Papel} & \textbf{Qué inyecta} \\
\midrule
\texttt{frio} & control base & Nada. \\
\midrule
\texttt{nearest} & \textbf{experimental} & Las 5 configuraciones de los vecinos más
cercanos. \textbf{Idéntica a la del experimento 1}, sin un solo cambio. \\
\midrule
\texttt{frequent} & control & Las 5 ganadoras más frecuentes del pliegue, sin mirar
las meta-características del conjunto consultado. \\
\midrule
\texttt{random} & ablación & Las ganadoras de 5 conjuntos elegidos al azar del
pliegue, con semilla 202. Aísla la contribución de la similitud. \\
\midrule
\texttt{hybrid} & propuesta & Las 2 primeras de \texttt{frequent} seguidas de las de
\texttt{nearest} en orden de cercanía, sin repetir, hasta reunir 5 distintas
(Sección~\ref{sec:hibrido}). \\
\bottomrule
\end{tabularx}
\end{center}

\textbf{Que \texttt{nearest} no cambie es deliberado y es lo que permite atribuir.}
Si \texttt{nearest} contra \texttt{frio} mejora respecto del experimento 1, la causa
es el instrumento, porque el método es literalmente el mismo. Mezclar las dos
mejoras en una sola condición haría inatribuible cualquier resultado favorable.

\textbf{Que \texttt{frequent} siga siendo un control también es deliberado.} Ganó en
el experimento 1; retirarlo o ascenderlo ahora sería seleccionar sobre el resultado.

\subsection{Definición de \texttt{hybrid}}
\label{sec:hibrido}

Construcción determinista, sin ningún parámetro ajustable y sin ajuste alguno:

\begin{enumerate}
  \item Tómense las listas que ese pliegue produce para \texttt{frequent} ---ordenada
        por frecuencia descendente, con el desempate canónico--- y para
        \texttt{nearest} ---ordenada por distancia creciente---.
  \item Tómense las \textbf{2 primeras} de \texttt{frequent}.
  \item Añádanse las de \texttt{nearest} en orden, descartando las repetidas, hasta
        reunir 5 distintas.
  \item Si se agotan, continúese con el resto de \texttt{frequent}; si aun así
        faltaran, con las ganadoras restantes del pliegue en orden canónico.
\end{enumerate}

El reparto 2 y 3 se declara aquí y no se explora: la mayoría queda del lado de la
similitud porque es el método comprometido en el anteproyecto, y las 2 por frecuencia
son el mínimo que garantiza incluir las configuraciones más transferibles del
pliegue. \textbf{No se probará ningún otro reparto}; hacerlo y quedarse con el mejor
sería ajustar sobre el resultado.

% =========== 8. PLAN DE ANALISIS ============================
\section{Plan de análisis}
\label{sec:analisis}

\subsection{Unidad de observación}

Por cada conjunto de datos y condición: la \textbf{mediana de las diez semillas} de la
variable primaria. Una observación por conjunto y condición, \textbf{42 pares} por
contraste, frente a los 25 del experimento 1.

\subsection{Los dos contrastes confirmatorios}

Se declaran \textbf{dos}, y solo dos, con prueba de Wilcoxon pareada de una cola:

\begin{center}
\begin{tabularx}{\textwidth}{L{1.2cm} L{3.3cm} Y}
\toprule
\rowcolor{tablehead} & \textbf{Contraste} & \textbf{Qué pregunta} \\
\midrule
C1 & \texttt{nearest} contra \texttt{frio} &
¿Reduce el arranque en caliente meta-aprendido el costo de la búsqueda? Es la
hipótesis del trabajo, medida ahora con resolución. \\
\midrule
C2 & \texttt{nearest} contra \texttt{frequent} &
¿Aporta algo la \emph{similitud} por encima de un buen valor por defecto? Es la
pregunta que el experimento 1 dejó planteada, y el listón exigente. \\
\bottomrule
\end{tabularx}
\end{center}

\textbf{Corrección por multiplicidad:} Holm--Bonferroni sobre esa familia de dos, con
$\alpha = 0{,}05$. Se reportan los valores $p$ crudos y los ajustados, y también los
de dos colas, para que la lateralidad no sea un grado de libertad oculto.

\textbf{Tamaño del efecto:} correlación biserial por rangos para muestras pareadas,
reportada junto a cada contraste.

\subsection{Criterio de éxito declarado}

Un contraste confirmatorio se declara favorable si cumple \textbf{las dos}
condiciones:

\begin{enumerate}
  \item valor $p$ de una cola \textbf{ajustado por Holm} menor que $0{,}05$, y
  \item reducción de la \textbf{mediana del regret} de al menos el \textbf{30\,\%}.
\end{enumerate}

El 30\,\% se traslada del criterio del anteproyecto para que ambos experimentos se
juzguen con la misma exigencia de magnitud. Es una \textbf{decisión declarada}, no una
derivación: el regret y los ensayos son escalas distintas y un 30\,\% en una no
equivale a un 30\,\% en la otra. Se fija ahora precisamente por eso.

\subsection{Análisis secundarios, declarados como tales}

No sustituyen a los confirmatorios y no llevan corrección por multiplicidad, porque no
se usan para decidir sobre la hipótesis:

\begin{itemize}
  \item \textbf{Los mismos contrastes sobre los ensayos hasta el objetivo}, con la
        definición y el criterio del anteproyecto. Es la comparación directa con el
        experimento 1.
  \item \texttt{hybrid} contra \texttt{frequent} y \texttt{hybrid} contra
        \texttt{nearest}: si la propuesta aporta. Es el secundario principal.
  \item \texttt{random} contra \texttt{frio}: la ablación que aísla la similitud.
  \item \textbf{Desempeño a presupuesto reducido}: mejor valor alcanzado a los 5, 15 y
        30 ensayos. El de 30 permite comparar con el experimento 1 a igual presupuesto.
  \item \textbf{Robustez de la censura}: la variable secundaria recalculada con 120 y
        600. El regret no se censura, de modo que el análisis primario es inmune a
        esta regla.
\end{itemize}

\subsection{Covariables declaradas}

Ambas se miden \textbf{antes} de ejecutar:

\begin{itemize}
  \item \textbf{Anchura de la meseta al 99\,\%}: cuántas de las 144 configuraciones
        alcanzan el objetivo. Ojo: la meseta que el experimento 1 declaró estaba
        medida \textbf{al 95\,\%}, así que hay que recalcularla; usar aquella sería
        usar una covariable que ya no corresponde al objetivo vigente.
  \item \textbf{Fracción de ensayos inevaluables} por corrida.
\end{itemize}

\subsection{Análisis exploratorios, no confirmatorios}

Se reportarán marcados como exploratorios, y de ellos no se deriva ninguna
afirmación sobre la hipótesis: la ponderación o selección de meta-características
para mitigar la concentración de distancias, y la relación entre la anchura de la
meseta y el beneficio observado.

% =========== 9. RESULTADOS POSIBLES =========================
\section{Resultados posibles, declarados de antemano}
\label{sec:desenlaces}

\begin{center}
\begin{tabularx}{\textwidth}{L{4.6cm} Y}
\toprule
\rowcolor{tablehead} \textbf{Desenlace} & \textbf{Cómo se reporta} \\
\midrule
C1 y C2 favorables &
El arranque en caliente reduce el costo \textbf{y} la similitud aporta por encima de
un buen valor por defecto. Es el desenlace más fuerte. \\
\midrule
C1 favorable, C2 no &
El arranque en caliente ayuda, pero \textbf{no por la similitud}: un valor por
defecto obtendría lo mismo. Refuerza el hallazgo del experimento 1 y es un
resultado sustantivo sobre los límites del meta-aprendizaje en este espacio. \\
\midrule
Ninguno favorable &
\textbf{Resultado negativo}, y se reporta como tal, con la misma extensión que uno
favorable. Con el instrumento ya corregido, pasaría a ser evidencia de que el efecto
no está ahí, y no de que no se podía medir. \\
\midrule
C2 favorable, C1 no &
Improbable y se reportaría con cautela: significaría que la similitud supera al valor
por defecto sin que ninguno de los dos supere al arranque en frío. \\
\bottomrule
\end{tabularx}
\end{center}

Ningún desenlace autoriza a modificar este documento después de conocerlo.

% =========== 10. COSTO ======================================
\section{Costo previsto y su contingencia}
\label{sec:costo}

El experimento 2 son $42 \times 10 \times 5 \times 60 = 126\,000$ ensayos, más
$42 \times 10 \times 144 = 60\,480$ evaluaciones de las rejillas de referencia. Con
el rendimiento medido en el experimento 1 ---15\,000 ensayos en 13,4 minutos--- la
estimación es de unas tres a cuatro horas.

El costo se \textbf{mide antes de lanzar}, cronometrando una condición sobre los
conjuntos más caros, como se hizo en el experimento 1.

\textbf{Contingencia declarada:} si la proyección superara las \textbf{12 horas}, las
semillas bajan de diez a cinco ---las de la v1.0--- y todo lo demás se mantiene. Esta
decisión se toma \textbf{solo sobre el costo medido}, jamás sobre resultado alguno, y
por eso se fija ahora.

% =========== 11. FUERA DE ALCANCE ===========================
\section{Qué queda fuera de este documento}
\label{sec:fuera}

\begin{itemize}
  \item \textbf{El experimento 1.} Lo gobierna el pre-registro v1.0 y no se toca.
  \item \textbf{Un espacio de búsqueda nuevo.} Se evaluó construir un
        \texttt{tabular-v2} con hiperparámetros continuos y se descartó por las
        razones del ADR-004. \texttt{tabular-v1} sigue vigente.
  \item \textbf{El caso de anomalías en imágenes} de la semana 11, que verifica la
        generalización por tarea y no es un contraste de esta hipótesis.
  \item \textbf{Cualquier ponderación de meta-características} como condición
        confirmatoria: no encontré un criterio de ajuste fijable \emph{a priori} sin
        arbitrariedad, y se relega a exploratorio.
\end{itemize}

% =========== 12. HISTORIAL ==================================
\section{Historial de versiones}
\label{sec:historial}

\begin{center}
\begin{tabularx}{\textwidth}{L{1.5cm} L{3.0cm} Y}
\toprule
\rowcolor{tablehead} \textbf{Versión} & \textbf{Cierre} & \textbf{Qué fija} \\
\midrule
\preregversion & \preregfecha &
Primera y única versión de este documento. Fija el protocolo del experimento 2:
objetivo al 99\,\% del óptimo exhaustivo, presupuesto de 60 ensayos, \emph{regret}
normalizado como variable primaria, corpus de 42, diez semillas, cinco condiciones y
dos contrastes confirmatorios con corrección de Holm. \\
\bottomrule
\end{tabularx}
\end{center}

\vspace{4pt}
{\small\color{gris} Los pre-registros v0.1, v0.2 y v1.0, que gobiernan el experimento
1 y todo lo anterior a él, viven en \texttt{Preregistro.tex} con sus propias
etiquetas. Este documento no los sustituye ni los enmienda.\par}

\end{document}
